% PATCH: C-84
% Reviewer: — (final round-2 corrections; brief intake/pending/final_corrections_round2_2026-11-04.md)
% Target section: sections/introduction.tex, sections/conclusions.tex, sections/preamble.tex
%   (closing.tex needed no change)
% Status: applied
%
% Applied, by audit ID (details under "### <ID>" in intake/pending/final_audit_r2_*):
% - B.2 (S-06): Introduction and Conclusions no longer say the gait is generated by CPGs: the walking
%   gait is a fixed open-loop crawl pattern (as Methodology says); the "BG above, CPG below" stack
%   sentence removed (its citations remain cited in Table 1). Also fixes S-40 #09 = L-37.
% - N-1: the MLP classifies terrain and feeds the gait-decision module; it does not compute BG salience.
% - B.16 (S-38 = M-19): contribution (iii) is "autonomous mode selection" (no "energy-aware");
%   B-15: (i)-(iii) use the same separator.
% - D-18 (Diego): Conclusions claims without a measurement made qualitative ("minimal deviations",
%   "consistent safety distance", "significantly reduced", "fewer undesired oscillations" removed),
%   anchored to figures/sections.
% - B.15 (S-25): robustness sentence rewritten with the brief's wording (combined perturbation, not per
%   sensor).
% - B.17a (S-33 = M-18 = L-13): CRediT names and order copied from the title page.
% - B.17b (B-13): AI disclosure names Claude (Anthropic) and its coding use; brief's proposed text.
% - D-47 (Diego): GPR expanded. L-10: STN/STR expanded in the abstract (Highlights unchanged: 85-char
%   limit).
% - S-40 #01 = L-44: final gerund of the Conclusions split into its own sentence.
% - L-20: "(Section~\ref{sec:results})" -> ssec:QuantReal where the energy result is meant.
% - B-4 = R-15, L-27, L-15, L-19, L-28 = M-21, L-31, B-14, B-10, S-37: second "developed at ETH Zurich"
%   removed, "ANYmal-on-wheels", MPC defined at first use, 180$^\circ$, "stimuli of different types",
%   "civil structures", "This model and its computational formulation", "made on", "fluctuating", en
%   dash, double space.
%
% Not applied (flag in PR): L-31 "Horizontal comparison" in the Table 1 caption -- that is Reviewer 3's
%   own wording in round 1 ("horizontal classification and comparative tables").
%
% Unified diff (sections/ only):
% diff --git a/sections/conclusions.tex b/sections/conclusions.tex
% index b63190b..0d27cfd 100644
% --- a/sections/conclusions.tex
% +++ b/sections/conclusions.tex
% @@ -1,21 +1,21 @@
%  \section{Conclusions}
%  \label{sec:conclusions}
%  
% -The experimental results, both in simulation and on the constructed platform, demonstrate that the proposed neural architecture enables the quadruped robot to perform stable and adaptive navigation in controlled environments with different stimulus configurations and topographical variations. In scenarios featuring stimuli of different natures, the physical robot reliably reaches the appetitive stimulus while avoiding obstacles and evading the aversive stimulus, maintaining efficient trajectories and a constant frontal orientation toward the target. Overall, these results confirm that combining bio-inspired control with multimodal sensing constitutes an effective strategy for autonomous locomotion in robotic inspection contexts. 
% +The experimental results, both in simulation and on the constructed platform, demonstrate that the proposed neural architecture enables the quadruped robot to perform stable and adaptive navigation in controlled environments with different stimulus configurations and topographical variations. In scenarios featuring stimuli of different types, the physical robot reliably reaches the appetitive stimulus while avoiding obstacles and evading the aversive stimulus, maintaining efficient trajectories and a constant frontal orientation toward the target. Overall, these results confirm that combining bio-inspired control with multimodal sensing constitutes an effective strategy for autonomous locomotion in robotic inspection contexts. 
%  
% -The observed performance highlights the effectiveness of the obstacle-perception module and the locomotion scheme controlled by central neural networks. \revblue{Coupled to the locomotion module, the LiDAR-driven neuronal ring facilitated trajectories with minimal deviations from the optimal path, maintaining a consistent safety distance from non-relevant obstacles and a persistent orientation toward the target without oscillations or unnecessary course corrections (neuron-level traces in Section~\ref{sec:results}).} This coherence between perception and motor control confirms that integrating a directional sensory system with a central pattern generator modulated by stimulus type is sufficient to ensure stable and efficient displacement even in the absence of environmental conflict. 
% +The observed performance highlights the effectiveness of the obstacle-perception module and the locomotion scheme controlled by central neural networks. \revblue{Coupled to the locomotion module, the LiDAR-driven neural ring steered the robot around task-irrelevant obstacles while keeping its orientation toward the target, as the trajectories in Figures~\ref{fig:RES_EST_MUL} and~\ref{fig:RESPUESTAS_EST_UNI_real} show. This coherence between perception and motor control indicates that integrating a directional sensory system with a stimulus-modulated locomotion network supports stable displacement in the absence of environmental conflict.} 
%  
%  In scenarios presenting simultaneous appetitive and aversive stimuli, the basal-ganglia-inspired winner-take-all arbitration module (Section~\ref{sssec:Basal_Ganglia_Module}) consistently prioritized the most relevant objective, its predominant channel activation inhibiting competing channels and reducing the probability of erratic trajectory changes; \revblue{in simulation, the full circuit resolved such conflicts with a mean decision latency of $75.9$~ms (95\%~CI $[59.3, 92.6]$), while on the physical robot the mean neural latency of terrain-driven mode transitions was 184~ms (quadrupedal$\rightarrow$rolling) and 534~ms (rolling$\rightarrow$quadrupedal) (Figure~\ref{fig:FourPanelLatencies}).} The alignment between decision-neuron activity peaks and executed maneuvers supports the hypothesis that a basal-ganglia-inspired architecture combined with multimodal sensory processing offers clear advantages in environments where goal prioritization is critical for mission success. 
%  
% -On irregular surfaces and steep slopes, the system exhibited an adaptive response grounded in the direct integration of inertial data into the locomotion control module. \revblue{Continuous monitoring of roll, pitch, and the standard deviation of vertical acceleration enabled early detection of instability conditions.} \revblue{When a threshold was exceeded, the neural architecture switched gait: high $\sigma_{az}$ on rugged terrain engaged the quadrupedal mode, while sustained tilt on the slope engaged differential rolling. Each switch showed a gradual activation of the modulator neurons ($X_{0}$–$X_{2}$) and a controlled inhibition of the previous gait. These transitions} not only maintained postural stability but also significantly reduced lateral and longitudinal oscillation amplitudes, minimizing the risk of loss of balance. The relationship between IMU signal intensity and postural-response activation confirms that combining kinematic and dynamic criteria is an effective strategy for locomotor adaptation in adverse terrains, ensuring both structural integrity of the robot and continuity of the mission. Likewise, it is noted that the first version of the gait-decision network implemented in simulation produces variant transient states during transitions to and from the quadrupedal mode; therefore, exploring additional IMU-derived variables or signal-processing approaches could yield more robust terrain-state inference. 
% +On irregular surfaces and steep slopes, the system exhibited an adaptive response grounded in the direct integration of inertial data into the locomotion control module. \revblue{Continuous monitoring of roll, pitch, and the standard deviation of vertical acceleration enabled early detection of instability conditions.} \revblue{When a threshold was exceeded, the neural architecture switched gait: high $\sigma_{az}$ on rugged terrain engaged the quadrupedal mode, while sustained tilt on the slope engaged differential rolling. Each switch showed a gradual activation of the modulator neurons ($X_{0}$--$X_{2}$) and a controlled inhibition of the previous gait. These transitions} \revblue{maintained postural stability in the reported trials, without abrupt roll or pitch excursions at the switching instant (Section~\ref{sssec:transition_stability})}. The relationship between IMU signal intensity and postural-response activation confirms that combining kinematic and dynamic criteria is an effective strategy for locomotor adaptation in adverse terrains, ensuring both structural integrity of the robot and continuity of the mission. Likewise, the first version of the gait-decision network implemented in simulation produces fluctuating transient states during transitions to and from the quadrupedal mode; therefore, exploring additional IMU-derived variables or signal-processing approaches could yield more robust terrain-state inference. 
%  
% -These results are consistent with those reported by \citet{sarvestani2013computational}, who demonstrated that a basal-ganglia-based arbitration scheme can manage sensory conflicts and robustly prioritize motor behaviors. \revblue{However, this work contributes three differentiating properties beyond the state of the art. First, biological interpretability: the basal-ganglia-inspired arbitration and the locomotor network modulated by the Naka--Rushton function keep each decision traceable to named neural units, and the latter yields smoother transitions with fewer undesired oscillations. Second, modularity: perception, arbitration, and motor generation are decoupled, so a single stage (for example, the inertial stability monitor that enables early detection of disturbances) can be upgraded without redesigning the others. Third, autonomous mode selection that exploits the efficiency of wheeled locomotion on flat surfaces while preserving legged capability on complex terrain, switching gait from inertial and visual cues, with the energy cost of each mode transition measured on the physical platform (Section~\ref{sec:results}). Together, these properties provide a computational foundation for autonomous gait arbitration in hybrid biomorphic platforms.} 
% +These results are consistent with those reported by \citet{sarvestani2013computational}, who demonstrated that a basal-ganglia-based arbitration scheme can manage sensory conflicts and robustly prioritize motor behaviors. \revblue{However, this work contributes three differentiating properties beyond the state of the art. First, biological interpretability: the basal-ganglia-inspired arbitration and the locomotor network modulated by the Naka--Rushton function keep each decision traceable to named neural units. Second, modularity: perception, arbitration, and motor generation are decoupled, so a single stage (for example, the inertial stability monitor that enables early detection of disturbances) can be upgraded without redesigning the others. Third, autonomous mode selection that exploits the efficiency of wheeled locomotion on flat surfaces while preserving legged capability on complex terrain, switching gait from inertial and visual cues, with the energy cost of each mode transition measured on the physical platform (Section~\ref{ssec:QuantReal}). Together, these properties provide a computational foundation for autonomous gait arbitration in hybrid biomorphic platforms.} 
%  
%  Section~\ref{ssec:Limitations} details the present work's limitations, spanning the physical platform's hardware envelope, the scope of experimental validation, manual parameter tuning, and the absence of a formal sensor-fusion framework; the future-work priorities they motivate are consolidated below.
%  
% -\revblue{The findings show that lightweight, bio-inspired neural networks that arbitrate between locomotion modes and reactive behaviors from visual and inertial cues can manage multimodal gait transitions on a real robotic platform, with the neural circuits executed on an external computer.} The perception-decision pipeline additionally retained task-level performance under combined IMU-drift and LiDAR-noise perturbation, and, conditionally on the scenario, camera illumination distortion (Section~\ref{ssec:NoiseRobustnessEval}), evaluated jointly rather than as isolated single-sensor degradations. \revblue{Although the physical prototype displays hardware-level boundaries (Section~\ref{ssec:Limitations}),} these results substantiate that complex, multi-stimulus mobility decisions can be \revblue{taken for such resource-constrained platforms; onboard feasibility was assessed only for the MLP terrain classifier on the ESP32-S3 (Section~\ref{sssec:mlp_eval}), not for the full neural pipeline}.
% +\revblue{The findings show that lightweight, bio-inspired neural networks that arbitrate between locomotion modes and reactive behaviors from visual and inertial cues can manage multimodal gait transitions on a real robotic platform, with the neural circuits executed on an external computer.} \revblue{The perception-decision pipeline additionally retained task-level performance under the combined perturbation in most scenarios (failures concentrated in Appetitive timeouts and Rugged Terrain tipovers), evaluated jointly rather than per sensor (Section~\ref{ssec:NoiseRobustnessEval}).} \revblue{Although the physical prototype displays hardware-level boundaries (Section~\ref{ssec:Limitations}),} these results substantiate that complex, multi-stimulus mobility decisions can be \revblue{made on such resource-constrained platforms; onboard feasibility was assessed only for the MLP terrain classifier on the ESP32-S3 (Section~\ref{sssec:mlp_eval}), not for the full neural pipeline}.
%  
% -The modular nature of the architecture facilitates its adaptation to different robotic platforms and sensor types, broadening its applicability across mission scenarios; the current pipeline of inertial and vision data (including stagnation detection) already allows new sensor modalities to be incorporated without modifying the system's core arbitration and sensorimotor fusion principles. Beyond mobility, the framework is extensible toward inspection-specific sensing: incorporating NDT modalities (infrared, ultrasound, hyperspectral) and fusing them over 2.5D maps will enable more reliable damage diagnostics. Future-work priorities include: long-term field validation under unconstrained, real-world outdoor conditions; migrating the arbitration logic to industrial-grade platforms with closed-loop dynamic gait replanning; integrating a formal sensor-fusion framework; exploring energy-efficient lateral-locomotion alternatives that reduce the Omnidirectional mode's parasitic-friction overhead (Section~\ref{sec:results}); and adaptive threshold learning and policy optimization to reduce manual tuning, together with safety and traceability protocols aligned with current standards, evolving the platform from a functional prototype into a practical autonomous inspection tool.
% +The modular nature of the architecture facilitates its adaptation to different robotic platforms and sensor types, broadening its applicability across mission scenarios; the current pipeline of inertial and vision data (including stagnation detection) already allows new sensor modalities to be incorporated without modifying the system's core arbitration and sensorimotor fusion principles. Beyond mobility, the framework is extensible toward inspection-specific sensing: incorporating NDT modalities (infrared, ultrasound, hyperspectral) and fusing them over 2.5D maps will enable more reliable damage diagnostics. Future-work priorities include: long-term field validation under unconstrained, real-world outdoor conditions; migrating the arbitration logic to industrial-grade platforms with closed-loop dynamic gait replanning; integrating a formal sensor-fusion framework; exploring energy-efficient lateral-locomotion alternatives that reduce the Omnidirectional mode's parasitic-friction overhead (Section~\ref{ssec:QuantReal}); and adaptive threshold learning and policy optimization to reduce manual tuning, together with safety and traceability protocols aligned with current standards. \revblue{Together, these steps would evolve the platform from a functional prototype into a practical autonomous inspection tool.}
%  
%  \section*{Acknowledgements}
%  
% @@ -23,12 +23,12 @@ The authors would like to thank the Neurocontrol Research Group at the Universid
%  
%  \section*{Author Contributions}
%  
% -\textbf{Samuel F. Carlos Rodriguez:} Methodology, Software, Investigation, Writing--original draft, Supervision.
% -\textbf{Diego A. Caro Vaca:} Software, Investigation, Writing--original draft, Supervision.
% -\textbf{Sebastian Grajales Pulgarín:} Methodology, Software, Investigation, Resources.
% -\textbf{David Caicedo Samboni:} Investigation, Formal analysis, Writing--original draft, Visualization.
% +\textbf{Samuel Federico Carlos-Rodríguez:} Methodology, Software, Investigation, Writing--original draft, Supervision.
% +\textbf{Diego Alejandro Caro-Vaca:} Software, Investigation, Writing--original draft, Supervision.
% +\textbf{David Caicedo-Samboní:} Investigation, Formal analysis, Writing--original draft, Visualization.
% +\textbf{Sebastián Grajales-Pulgarín:} Methodology, Software, Investigation, Resources.
%  \textbf{Julián Hurtado-López:} Supervision, Project administration.
% -\textbf{David F. Ramirez-Moreno:} Project administration, Resources.
% +\textbf{David Fernando Ramírez-Moreno:} Project administration, Resources.
%  \textbf{All authors:} Conceptualization, Writing--review \& editing, Validation.
%  
%  \section*{Declaration of Competing Interest}
% @@ -37,7 +37,7 @@ The authors declare that they have no known competing financial interests or per
%  
%  \section*{AI-assisted writing disclosure}
%  
% -During the writing of this manuscript and its subsequent revisions, we used artificial intelligence-based tools, including ChatGPT (OpenAI), as editorial support to improve the grammar, clarity, tone, coherence, and overall writing quality of the text. These tools were used solely for linguistic and stylistic refinement and were not used to generate scientific results, data, analyses, or conclusions. The authors critically reviewed, edited, and approved the final content and take full responsibility for the accuracy and integrity of the manuscript.
% +\revblue{During the writing of this manuscript and its revisions, we used AI-based tools, including ChatGPT (OpenAI) and Claude (Anthropic), as editorial support to improve grammar, clarity, tone and coherence, and Claude (Anthropic) as a coding assistant to maintain the manuscript source and to run the analysis and plotting scripts that regenerate figures and tables from the experimental data. These tools were not used to design the experiments, collect the data, or draw the scientific conclusions. The authors critically reviewed, edited and approved all content, verified the outputs of the scripts, and take full responsibility for the accuracy and integrity of the manuscript.}
%  
%  %% For citations use: 
%  %%       \citet{<label>} ==> Lamport [21]
% diff --git a/sections/introduction.tex b/sections/introduction.tex
% index a868b83..ae8ebf9 100644
% --- a/sections/introduction.tex
% +++ b/sections/introduction.tex
% @@ -1,11 +1,11 @@
%  \section{Introduction}
%  \label{sec:introduction}
%  
% -Aging infrastructure worldwide, including bridges, industrial plants, tunnels, and civil constructions, requires frequent inspection and maintenance to prevent catastrophic failures. The manual execution of these tasks exposes human workers to hazardous conditions, incurs high operational costs, and is limited in coverage and repeatability \citep{Halder2023Construction,Wang2025Autonomous}. Autonomous robotic inspection has therefore emerged as a critical technology, offering scalability, continuous operation, and the capacity to acquire high-precision data in environments that are inaccessible or dangerous for humans \citep{Halder2025Crack}. The integration of autonomous sensing platforms into non-destructive testing (NDT) workflows has been identified as a key enabler for next-generation infrastructure monitoring, combining visual detection, inertial navigation, and point-cloud mapping into a unified inspection pipeline.
% +Aging infrastructure worldwide, including bridges, industrial plants, tunnels, and civil structures, requires frequent inspection and maintenance to prevent catastrophic failures. The manual execution of these tasks exposes human workers to hazardous conditions, incurs high operational costs, and is limited in coverage and repeatability \citep{Halder2023Construction,Wang2025Autonomous}. Autonomous robotic inspection has therefore emerged as a critical technology, offering scalability, continuous operation, and the capacity to acquire high-precision data in environments that are inaccessible or dangerous for humans \citep{Halder2025Crack}. The integration of autonomous sensing platforms into non-destructive testing (NDT) workflows has been identified as a key enabler for next-generation infrastructure monitoring, combining visual detection, inertial navigation, and point-cloud mapping into a unified inspection pipeline.
%  
%  Among existing robotic platforms, legged robots (and quadrupeds in particular) have proven the most effective for semi-structured environments that combine flat corridors with obstacles, stairs, and uneven ground. ANYmal, developed at ETH Z\"{u}rich, was specifically designed for harsh industrial environments and has demonstrated reliable operation in oil and gas plants and search-and-rescue scenarios \citep{Hutter2017ANYmal}. More recently, Team CERBERUS deployed a multi-agent system of ANYmal robots augmented with LiDAR, cameras, and IMU to win the DARPA Subterranean Challenge, establishing the quadruped as the reference platform for autonomous exploration in confined and semi-structured spaces \citep{Tranzatto2022CERBERUS}. Despite these advances, quadrupedal robots face a fundamental efficiency limitation: walking gaits are energetically costly on flat surfaces, which restricts their operational autonomy for tasks that combine open traversal (e.g., factory floors) with the negotiation of obstacles or slopes \citep{Vogel2024Ladder}. This trade-off between terrain adaptability and locomotion efficiency motivates the development of hybrid platforms capable of switching between legged and wheeled modes according to terrain properties.
%  
% -Wheeled-legged robots address this limitation by integrating actuated limbs (providing the dexterity required to negotiate irregular or sloped terrain) with wheels that constitute the most efficient effector on flat surfaces \citep{Bjelonic2023Survey}. The ANYmal-on-wheels family, developed at ETH Z\"{u}rich, is the most representative line in this category. \citet{Bjelonic2019KeepRollin} introduced the first whole-body controller capable of simultaneously combining walking and rolling with non-steerable torque-controlled wheels; subsequent work \citep{bjelonic2019} extended this to online trajectory optimization for hybrid maneuvers validated in the DARPA SubT, and \citet{Bjelonic2022OfflineMotion} later incorporated offline motion libraries and MPC to achieve complex skills such as stair climbing and dynamic 180° turns. The most recent and impactful result in this family, \citet{Lee2024Science}, demonstrated a complete pipeline of adaptive locomotion control, mobility-aware local planning, and city-scale global navigation via deep reinforcement learning (RL) with privileged learning, achieving smooth walking-to-driving transitions as an emergent property of a single learned policy. The CENTAURO platform from the Italian Institute of Technology \citep{Kashiri2019CENTAURO,Klamt2019CENTAURO}, a centaur-type quadruped with steerable omnidirectional wheels and a high-payload manipulator arm, offers a complementary design point with explicit disaster-response validation and provides the best published quantitative evidence for the energy/terrain trade-off that motivates a hybrid robot. At the compact end of the spectrum, Ascento \citep{Klemm2019Ascento} is a bipedal wheeled-legged platform for indoor inspection, demonstrating that the design space spans from bipedal (Ascento) to centaur (CENTAURO) and quadrupedal (ANYmal-wheels) morphologies. Recent platforms such as those published by \citet{Liu2024Max},  \citet{Cui2025FLORES}, and \citet{Jiang2025RTaichi} further expand the design space with alternative articulation schemes and RL-based controllers, confirming that the field is converging on hybrid locomotion as the dominant paradigm for versatile robot operation.
% +Wheeled-legged robots address this limitation by integrating actuated limbs (providing the dexterity required to negotiate irregular or sloped terrain) with wheels that constitute the most efficient effector on flat surfaces \citep{Bjelonic2023Survey}. The ANYmal-on-wheels family is the most representative line in this category. \citet{Bjelonic2019KeepRollin} introduced the first whole-body controller capable of simultaneously combining walking and rolling with non-steerable torque-controlled wheels; subsequent work \citep{bjelonic2019} extended this to online trajectory optimization for hybrid maneuvers validated in the DARPA SubT, and \citet{Bjelonic2022OfflineMotion} later incorporated offline motion libraries and model predictive control (MPC) to achieve complex skills such as stair climbing and dynamic $180^{\circ}$ turns. The most recent and impactful result in this family, \citet{Lee2024Science}, demonstrated a complete pipeline of adaptive locomotion control, mobility-aware local planning, and city-scale global navigation via deep reinforcement learning (RL) with privileged learning, achieving smooth walking-to-driving transitions as an emergent property of a single learned policy. The CENTAURO platform from the Italian Institute of Technology \citep{Kashiri2019CENTAURO,Klamt2019CENTAURO}, a centaur-type quadruped with steerable omnidirectional wheels and a high-payload manipulator arm, offers a complementary design point with explicit disaster-response validation and provides the best published quantitative evidence for the energy/terrain trade-off that motivates a hybrid robot. At the compact end of the spectrum, Ascento \citep{Klemm2019Ascento} is a bipedal wheeled-legged platform for indoor inspection, demonstrating that the design space spans from bipedal (Ascento) to centaur (CENTAURO) and quadrupedal (ANYmal-on-wheels) morphologies. Recent platforms such as those published by \citet{Liu2024Max}, \citet{Cui2025FLORES}, and \citet{Jiang2025RTaichi} further expand the design space with alternative articulation schemes and RL-based controllers, confirming that the field is converging on hybrid locomotion as the dominant paradigm for versatile robot operation.
%  
%  A critical observation emerges from this body of work: in all published wheeled-legged systems, locomotion mode switching is resolved either heuristically through hand-coded threshold logic \citep{Kashiri2019CENTAURO}, by model predictive control \citep{bjelonic2019,Bjelonic2022OfflineMotion}, or by end-to-end RL policies that encode switching implicitly as a black-box behavior \citep{Lee2024Science,Humphreys2025BioInspired}. None of these approaches employs a neurobiologically interpretable arbitration architecture. \revblue{Table~\ref{table:action_selection_comparison} contrasts these approaches with central pattern generators (CPGs), including feedback-modulated variants that trigger gait transitions from vestibular or load feedback, and with the basal-ganglia winner-take-all arbitration adopted here.}
%  
% @@ -29,9 +29,9 @@ Basal ganglia (WTA disinhibition), \emph{this work} & Salience-driven winner-tak
%  
%  The basal ganglia, a subcortical neural circuit present in all vertebrates, are the canonical neural substrate for action selection: through a salience-driven winner-take-all (WTA) mechanism implemented via disinhibition, they select and sustain a single behavioral channel from a set of competing alternatives while suppressing all others. To the best of the authors' knowledge, no prior work has applied a basal-ganglia-type architecture to arbitrate among rolling, walking, and omnidirectional locomotion modes in a hybrid quadruped. This constitutes the principal novelty of the present contribution.
%  
% -Neural control architectures seek to emulate the information-processing structures found in biological systems in order to generate adaptive and robust motor responses \citep{mead1990neuromorphic,ijspeert2008}. The GPR computational model of the basal ganglia (BG), proposed by \citet{Gurney2001a,Gurney2001b}, reformulates action selection as the resolution of competition among salience-coded channels through a feedforward off-center/on-surround network with selection and control pathways. \citet{Prescott2006Robot} provided the first embodied validation of the GPR model, implementing it in a mobile robot performing a foraging task with five competing behaviors, demonstrating clean switching, appropriate behavioral persistence, and conflict-resolution dynamics analogous to those observed in animals. \citet{Girard2003Basal} compared BG-based action selection against a classical WTA and showed that the BG's thalamocortical loops confer specific adaptive properties such as avoidance of dithering and energetic savings that a simple WTA cannot reproduce. Building upon this theoretical foundation, \citet{sarvestani2013computational} developed a computational network model that operationalizes these selection mechanics. This distinction and its subsequent computational formulation are directly relevant to the present work, as they provide the structural blueprint for our real-time mode arbitration. More recent updates to the GPR paradigm include a biologically parsimonious three-pathway model (Go/NoGo/hyperdirect) \citep{Baston2015Biologically} and a neurorobotic study of simulated dopamine modulation that shows how tonic dopamine levels control the speed and frequency of mode switching \citep{Prescott2024Simulated}, opening a pathway to meta-control of switching aggressiveness in future work. For the rhythmic generation of the quadrupedal walking gait, the present platform integrates central pattern generators (CPGs), neural circuits that produce coordinated rhythmic motor patterns from simple inputs \citep{ijspeert2008}; while feedback-modulated CPGs can themselves arbitrate discrete gait transitions \citep{Fukui2019VestibularCPG,Owaki2017SpontaneousGait}, the present design deliberately keeps that decision with the basal-ganglia stage instead, so that mode arbitration remains a single, biologically interpretable locus rather than an emergent property distributed across the oscillator network, a property advantageous for diagnosing and validating the physical mode transitions reported in the Results section. This combination of BG arbitration above and CPG execution below forms a hierarchical neural stack that is anatomically plausible and computationally modular \citep{Bellegarda2022CPGRL,Manoonpong2013Neural}.
% +Neural control architectures seek to emulate the information-processing structures found in biological systems in order to generate adaptive and robust motor responses \citep{mead1990neuromorphic,ijspeert2008}. The \revblue{Gurney--Prescott--Redgrave (GPR)} computational model of the basal ganglia (BG), proposed by \citet{Gurney2001a,Gurney2001b}, reformulates action selection as the resolution of competition among salience-coded channels through a feedforward off-center/on-surround network with selection and control pathways. \citet{Prescott2006Robot} provided the first embodied validation of the GPR model, implementing it in a mobile robot performing a foraging task with five competing behaviors, demonstrating clean switching, appropriate behavioral persistence, and conflict-resolution dynamics analogous to those observed in animals. \citet{Girard2003Basal} compared BG-based action selection against a classical WTA and showed that the BG's thalamocortical loops confer specific adaptive properties such as avoidance of dithering and energetic savings that a simple WTA cannot reproduce. Building upon this theoretical foundation, \citet{sarvestani2013computational} developed a computational network model that operationalizes these selection mechanics. This model and its computational formulation are directly relevant to the present work, as they provide the structural blueprint for our real-time mode arbitration. More recent updates to the GPR paradigm include a biologically parsimonious three-pathway model (Go/NoGo/hyperdirect) \citep{Baston2015Biologically} and a neurorobotic study of simulated dopamine modulation that shows how tonic dopamine levels control the speed and frequency of mode switching \citep{Prescott2024Simulated}, opening a pathway to meta-control of switching aggressiveness in future work. \revblue{The quadrupedal walking gait itself is executed by a fixed, open-loop crawl pattern (Section~\ref{ssec:ServoCompensation}), not by CPGs. Although feedback-modulated CPGs can arbitrate discrete gait transitions \citep{Fukui2019VestibularCPG,Owaki2017SpontaneousGait}, the present design keeps that decision in the basal-ganglia stage, so that mode arbitration remains a single, biologically interpretable locus rather than an emergent property of an oscillator network; a single locus is advantageous for diagnosing and validating the physical mode transitions reported in Section~\ref{sec:results}.}
%  
% -Against this background, the present work proposes a hybrid quadrupedal robotic platform with four three-degree-of-freedom limbs, each incorporating an omnidirectional wheel at the distal segment (tibia), enabling transitions among three discrete locomotion modes: quadrupedal walking (articulated mode), differential rolling, and omnidirectional motion (mobile mode). The platform integrates a sensory suite comprising an IMU, a LiDAR, and a camera for environmental perception. A multilayer perceptron (MLP) processes sensory data to detect terrain instability and compute per-channel salience inputs to the basal-ganglia network, which in turn executes mode selection via a WTA disinhibition mechanism. The complete system is validated in simulation (Gazebo) and on a physical prototype, demonstrating robust autonomous switching across flat surfaces, sloped terrain, and obstacle-laden environments. Relative to the current state of the art, the proposed architecture offers three differentiating properties: (i) \emph{biological interpretability}, following directly from the neuroscientific grounding established above; (ii) \emph{modularity}: perception, arbitration, and motor generation are decoupled, facilitating individual upgrade; and (iii) \emph{autonomous energy-aware mode selection}, leveraging the efficiency advantage of wheeled locomotion on flat surfaces while preserving legged capability for complex terrain, directly addressing the operational limitation documented in deployed inspection systems \citep{Hutter2017ANYmal,Tranzatto2022CERBERUS}.
% +Against this background, the present work proposes a hybrid quadrupedal robotic platform with four three-degree-of-freedom limbs, each incorporating an omnidirectional wheel at the distal segment (tibia), enabling transitions among three discrete locomotion modes: quadrupedal walking (articulated mode), differential rolling, and omnidirectional motion (mobile mode). The platform integrates a sensory suite comprising an IMU, a LiDAR, and a camera for environmental perception. \revblue{A multilayer perceptron (MLP) classifies terrain from inertial data on the physical platform and feeds the gait-decision module, while the basal-ganglia network arbitrates among the stimulus-driven behavioral channels via a WTA disinhibition mechanism.} The complete system is validated in simulation (Gazebo) and on a physical prototype, demonstrating robust autonomous switching across flat surfaces, sloped terrain, and obstacle-laden environments. Relative to the current state of the art, the proposed architecture offers three differentiating properties: (i) \emph{biological interpretability}, following directly from the neuroscientific grounding established above; (ii) \emph{modularity}, since perception, arbitration, and motor generation are decoupled, facilitating individual upgrade; and (iii) \revblue{\emph{autonomous mode selection}}, leveraging the efficiency advantage of wheeled locomotion on flat surfaces while preserving legged capability for complex terrain, directly addressing the operational limitation documented in deployed inspection systems \citep{Hutter2017ANYmal,Tranzatto2022CERBERUS}.
%  
%  The remainder of this article is organized as follows: Section~\ref{sec:methodology} details the methodology, including the physical robot design, the multi-body simulation setup, and the formal mathematical framework of the BG-inspired neural controller. Section~\ref{sec:results} presents the experimental results, separating virtual multi-body simulation tests from physical prototype validations. Finally, Section~\ref{sec:conclusions} outlines the conclusions and directions for future work.
%  
% diff --git a/sections/preamble.tex b/sections/preamble.tex
% index 863f673..16ec6cf 100644
% --- a/sections/preamble.tex
% +++ b/sections/preamble.tex
% @@ -82,7 +82,7 @@
%  
%  %% Abstract
%  \begin{abstract}
% -This work presents the design, simulation, and implementation of a bio-inspired quadruped robotic platform equipped with a hybrid wheel-leg locomotion system and controlled through a neural architecture modeled after basal ganglia mechanisms. The platform integrates sensory modules for environmental perception, arbitration modules for behavioral decision-making, and locomotor modules for gait generation, enabling adaptive responses to diverse stimuli and terrain conditions. Simulation and implementation demonstrate the robot's ability to switch autonomously between locomotion modes (quadrupedal walking, differential rolling, and omnidirectional motion) when facing obstacles, appetitive stimuli, or aversive conditions. \revblue{In simulation, an ablation study showed that removing the circuit's lateral inhibition or its STN/STR layers reduced success under stimulus conflict from 100\% to 25\% and 33\%, respectively. Against a priority-rule finite state machine (FSM) baseline under matched conditions, the neural arbitration succeeded in every three-stimulus conflict trial, where the FSM succeeded in none, and reached 85\% success versus 45\% in an inspection task with an initially hidden target.} Across simulated and real-robot trials, the controller \revblue{produced} terrain-specific, repeatable locomotion signatures, laying the groundwork for scaling toward robust real-world infrastructure inspection and operation in complex, unstructured environments.
% +This work presents the design, simulation, and implementation of a bio-inspired quadruped robotic platform equipped with a hybrid wheel-leg locomotion system and controlled through a neural architecture modeled after basal ganglia mechanisms. The platform integrates sensory modules for environmental perception, arbitration modules for behavioral decision-making, and locomotor modules for gait generation, enabling adaptive responses to diverse stimuli and terrain conditions. Simulation and implementation demonstrate the robot's ability to switch autonomously between locomotion modes (quadrupedal walking, differential rolling, and omnidirectional motion) when facing obstacles, appetitive stimuli, or aversive conditions. \revblue{In simulation, an ablation study showed that removing the circuit's lateral inhibition or its subthalamic/striatal (STN/STR) layers reduced success under stimulus conflict from 100\% to 25\% and 33\%, respectively. Against a priority-rule finite state machine (FSM) baseline under matched conditions, the neural arbitration succeeded in every three-stimulus conflict trial, where the FSM succeeded in none, and reached 85\% success versus 45\% in an inspection task with an initially hidden target.} Across simulated and real-robot trials, the controller \revblue{produced} terrain-specific, repeatable locomotion signatures, laying the groundwork for scaling toward robust real-world infrastructure inspection and operation in complex, unstructured environments.
%  \end{abstract}
%  
%  %%Graphical abstract
